\subsection{Morfologia do Rio Guaíba}
% Secretaria do Meio Ambiente e Infraestrutura: https://sema.rs.gov.br/g080-bh-guaiba

Localizada na região metropolitana do Rio Grande do Sul, a Bacia Hidrográfica do Guaíba constitui uma das Bacias mais
importantes do estado, possuíndo uma população de 1.3 milhões de pessoas nas cidades que percorrem o seu curso. Entre
elas se destacam os municípios da Região metropolitana de Porto Alegre (RMPOA), como Porto Alegre, Eldorado do Sul,
Guaíba e Canoas. 

\begin{figure}[htb]
	\caption{\label{guaiba_pol} Mapa Político do Guaíba}
	\begin{center}
	    \includegraphics[scale=0.48]{img/guaiba_pol.png}
	\end{center}
	\legend{Fonte: Instituto Brasileiro de Geografia e Estatística \cite{ibgeguai}}
\end{figure}


A mesma é principalmente formada pelas águas dos rios que desembocam no Guaíba, em especial aqueles que formam o Delta
do Jacuí, incluindo o Rio Caí, o Rio dos Sinos, o Rio Gravataí e o próprio Rio Jacuí. A região do Delta em si conta com
centros urbanos das cidades de Guaíba e Eldorado do Sul, além de contar com diversas ilhas, como a Ilha das Flores, a
Ilha Grande dos Marinheiros e a Ilha da Pintada. Alguns bairros de Porto Alegre também parte do sistema do ilhas, como o
Humaitá e o Navegantes.

% \begin{figure}[htb] \caption{\label{bacia_guaiba} Mapa da Bacia Hidrográfica do Lago Guaíba} \begin{center}
%   \includegraphics[scale=0.66]{img/bacia_guaiba.png} \end{center} \legend{Fonte: Secretaria do Meio Ambiente e
%   Infraestrutura do Rio Grande do Sul \cite{guaiba_bacia}} \end{figure}

Situa-se como destino das águas do Delta do Jacuí, destas, a descarga é composta primariamente pelo Rio Jacuí (84,6\%),
com o Rio dos Sinos, 7,5\%, Rio do Caí, com 5,2\%, e Rio Gravataí, com 2,7\% formando o restante da vazão. O Guaíba em
si pode ser tanto classificado como um Rio, quanto um lago. Historicamente, o mesmo já foi classificado de ambas as
formas, além de ria e estuário. O mesmo representa um ambiente de transição, com características de ambos rio e lago,
como um curso fluvial e uma dinâmica de armazenamento lacustre. Neste estudo, adota-se a nomenclatura "Rio Guaíba",
conforme padronização do IBGE em seus mapas oficiais \cite{ibgeguai}. Essa decisão é baseada na maior disponibilidade de
literatura relacionada à cheias fluvias e a utilização de fontes de dados provenientes do IBGE. Contudo, quando
apropriado, trabalhos envolvendo cheias em Lagos também serão utilizados de forma comparativa, como o de Güldal et al.
(2009) \cite{gudal}.

A Região Hidrográfica do Guaíba, da qual a Bacia Hidrográfica do Lago Guaíba faz parte, se extende pela parte
centro-leste do estado e inclui as Bacias Hidrográficas do Lago Guaíba, Sinos, Gravataí, Caí, Baixo-Jacuí, Alto-Jacuí,
Vacacaí, Pardo e Taquari-Antas, vide \autoref{regiao_guaiba}. A mesma constitui uma das duas Regiões Hidrográficas do
estado junto com a Região Hidrográfica do Rio Uruguai. 

Apresentando, em sua maioria, um padrão de drenagem dendrítico — caracterizado por uma rede de canais que se assemelha
aos ramos de uma árvore e que segue de forma adaptada ao relevo, vide \autoref{drain}.1 —, a Região Hidrográfica do
Guaíba concentra o escoamento das águas no Rio Guaíba, configurando um sistema de convergência típica em bacias com
vales em “V” e substrato predominantemente impermeável. Embora o padrão dendrítico seja predominante, variações locais
ocorrem devido à extensão e diversidade geomorfológica da região. A deposição de sedimentos, especialmente durante os
períodos de cheia, é marcante.

\begin{figure}[htb]
	\caption{\label{regiao_guaiba} Mapa da Região Hidrográfica do Guaíba}
	\begin{center}
	    \includegraphics[scale=0.35]{img/regiao_guaiba.png}
	\end{center}
	\legend{Fonte: Secretaria do Meio Ambiente e Infraestrutura do Rio Grande do Sul \cite{sema}}
\end{figure}

\subsection{Enchentes históricas}

A região da Bacia Hidrográfica do Guaíba é historicamente suscetível à cheias. A mais notável destas foi a Enchente de
1941 - marco hidrológico que permanece como referência catastrófica. Entre de Abril e Maio daquele ano, um período de
alta precipitação atingiu o estado e, somente na capital, a precipitação chegou a 791 milímetros. Combinado com o
escoamento das águas das chuvas ao redor do estado na Bacia Hidrográfica do Guaíba, o nível do mesmo chegou a 4,76
metros acima da cota normal \cite{dep}, submergindo bairros ribeirinhos como Centro Histórico, Menino Deus e Azenha por
22 dias consecutivos. Dos 272 mil habitantes então recenseados, 70 mil (25,7\%) foram desalojados, com transporte urbano
reduzido a embarcações de emergência. A contaminação do sistema de saneamento pela água pluvial desencadeou epidemias de
doenças hídricas nos meses subsequentes, enquanto empresas levaram meses para retomar operações. Este evento foi um dos
primeiros com bons registros na região. Subsequentemente, as cheias foram mais recorrentes: em 1967 o nível atingiu 3,13
m sem proteção do Muro da Mauá (então em construção); em 1984 registrou-se 2,50 m com ativação das comportas do sistema;
e em 2002 o pico de 2,46 m aproximou-se do limite operacional da infraestrutura.

A suscetibilidade regional é estrutural. Porto Alegre ocupa uma planície aluvial crítica com 35\% da área urbana abaixo
da cota 3 m e é localizada na confluência de quatro bacias hidrográficas (Jacuí, Gravataí, Sinos, Caí) que drenam 30\%
do território gaúcho \cite{dep_topografia}, além de ser cercada por arroios e pela Lagoa dos Patos ao sul. Este gargalo
hidrodinâmico coincide precisamente com o núcleo metropolitano, dado que na estreita região entre a Usina do Gasômetro e
a Ilha da Pintada, se encontra o maior ponto de represamento da região graças ao encontro dos rios. O crescimento urbano
e o desmatamento de encostas aumentaram a impermebilidade do solo e consequentemente diminuiram a capacidade de
inflitração urbana da região, aumentando o risco de uma área já vulnerável. Isto, somado à densidade populacional
presente nestas regiões, tornam crítico o monitoramento da cota das águas do Guaíba.

Nas últimas décadas, padrões pluviais intensificados pelas mudanças climáticas geraram eventos cada vez mais extremos:
em 2023 precipitações geradas por um ciclone extratropical levaram a dois eventos extremos no estado, um em Setembro e o
outro em Novembro. Afetando principalmente a Região do Vale do Taquari, na qual atingiu 359 mil pessoas, o excesso de
água causado pelas chuvas também chegou ao Rio Guaíba. No dia 21 de Novembro, a cota do rio chegou á 3,46 metros, o mais
alto desde a enchente de 1941. Porém, neste caso, os sistemas de prevenção de cheias da cidade consguiram parar grande
parte do avanço das águas.

\subsection{Sistemas de prevenção de cheias na RMPOA}
\lipsum[1-8]
% Casas de bombas, muro da Mauá, Galerias pluviais e como a falta de manutenção fez esses sistemas falharem -> Falta de
% Resiliência na malha urbana;

% Muro da Mauá. A enchente histórica trouxe à tona a ideia de construir um muro de proteção para Porto Alegre. Cerca de
% 30 anos depois do evento, entre 1971 e 1974, foi construído o Muro da Mauá, uma estrutura de concreto de 3 metros de
% altura e 2,6 km de extensão que faz parte de um sistema de contenção de águas.

% Diques. Atualmente, a cidade tem 24 km de diques que margeiam a FreeWay até a Avenida Castelo Branco e seguem pela
% Avenida Beira Rio; e mais 44 km de diques internos, às margens dos arroios que atravessam a cidade. São 68 km de
% diques, 14 comportas e 19 casas de bomba, de acordo com a prefeitura de Porto Alegre

\subsection{Enchentes de 2024}
\lipsum[1-6]

% INFORMATIVOS DO INMET -> https://portal.inmet.gov.br/informativos# (Tem que pesquisar no google cada um deles);

% Vale mensionar as enchentes de 2023 como introdução -> Talvez previsões de aumentos de enchentes na região sudeste do
% país? Tem um estudo que foi realizado no governo da Dilma;

% https://www.ncei.noaa.gov/monitoring-content/sotc/global/2024/apr/Brazil-Floods-RioGrandeSul-202404.pdf
% https://web.archive.org/web/20240930134441/https://www.nytimes.com/2024/05/02/world/americas/brazil-rain-floods.html
% https://www.bbc.com/news/world-latin-america-68935249
% https://www.reuters.com/graphics/BRAZIL-RAINS/FLOODS/egpbazlzbvq/

% Cronologia: https://www.bbc.com/portuguese/articles/cd1qwpg3z77o

% As enchentes em Porto Alegre exibiram características tanto de enchentes pluviais quanto fluviais, mas foi o último
% que gerou a maior parte da crise -> Foco no Guaíba. Isso se dever por ter tido a característica de uma enchente
% fluvial, isto é, o Rio foi aquilo que transbordou, mas também pode ser caracterizar como uma enchente pluvial urbana,
% no sentido de que o escoamento do Rio que levou ao aumento do Nível da Água e a falha do escoamento urbano que levou a
% água para dentro da cidade. 

% 5. Climate change, El Niño and infrastructure failures behind massive floods in southern Brazil